\section{Method}
\label{sec:method}

\paragraph{Paired updates and probes.}
Let $h^{(l)}(C)\in\mathbb R^{d}$ be the residual stream after block $l$ at the anchor for context $C$. For the increment $C_{k-1}\to C_k$ we form the paired update $\Delta h^{(l)}=h^{(l)}(C_k)-h^{(l)}(C_{k-1})$, which isolates what the added paragraph did to the model's state and removes the content that both contexts share. For each task and layer we fit an $\ell_2$-regularised logistic regression on standardised updates with class balancing. We probe every layer and report the one selected by grouped three-fold cross-validation on the training questions only; the depth profile is in \cref{app:readout}. The probe direction is the weight vector mapped back to activation space and normalised. We also compute the \emph{difference-of-means} direction $\vdm\propto\mathbb E[\Delta h\mid y{=}1]-\mathbb E[\Delta h\mid y{=}0]$ and, for comparison, probes on the raw states $h(C_{k-1})$ and $h(C_k)$ (\cref{app:rep}).

\paragraph{Nuisance analysis.}
A paired update also reflects surface properties of the increment: how much text was added, whether it mentions the answer, how lexically relevant it is to the question, where it was inserted, and how much the model's output distribution moved. We collect these into a \emph{nuisance vector} $z$ per increment (the full list is in \cref{app:probing}) and run two tests. A classifier that sees only $z$ measures how far surface properties alone go. The \emph{gain} is the AUROC added when the probe score, computed out of fold on training increments, is given to that classifier as an extra feature; it measures what the activations contribute beyond those properties.

\paragraph{Generalisation.}
Directions are learned on the training questions of 2-hop \musique{} and then applied, without retraining and at the layer chosen there, to other settings: a conversational protocol in which the model sees its previous answer before the new paragraph, \twowiki{}, 3--4-hop \musique{}, and the external benchmarks of \cref{sec:external}. Where the target has its own trajectories we also train probes on it, which gives the ceiling for transfer. Hidden sizes differ across model families, so there we replicate the whole procedure instead of transferring vectors.

\paragraph{Causal tests.}
A direction from which a property can be read is not necessarily one the model uses. We therefore ask three questions of a direction $v$, on 150 held-out questions, each with an \emph{insufficient} prompt that lacks one gold paragraph and the \emph{sufficient} prompt in which it has been added. \emph{Does pushing along $v$ change behaviour?} Steering adds $\alpha\,u\,v$ to the residual at the anchor and at every generated position, where $u$ is the gap between the sufficient and insufficient means along $v$, so that $\alpha=1$ moves the activation by one such gap; a random direction of the same norm is the control. \emph{Does the model need $v$?} Projection-out takes the sufficient prompt and sets its component along $v$ to the insufficient-state mean, from the chosen layer onward; if the model relied on $v$, it should start abstaining. \emph{Where in the update does the causal content lie?} Patching runs the insufficient prompt with its anchor residual replaced by the sufficient prompt's, either in full, or only in the component along $v$, or only in the component orthogonal to $v$; further variants restrict the replacement to a random subspace, to the top principal components of training updates, or to the unembedding direction of the gold answer's first token (\cref{app:probing}). Outcomes are the rate at which the model answers rather than saying INSUFFICIENT and the rate of correct answers. We test the logistic probe direction $\vlr$ and the difference-of-means direction $\vdm$ \cite{marks2023geometry,ActivationSteering,meng2022locating}.
